\subsection{Spectrum and Resolvent}\label{sec:spectrum-and-resolvent}

In this section we generalize the eigenvalue problems of linear algebra to operators on Banach spaces. One of our main goals is to generalize the following theorem.

\begin{Theorem}
\label{thm:symmetric-matrix-spectral-basis}
Let $A$ be an $n\times n$ symmetric matrix. Then $A$ has $n$ eigenvalues $\lambda_1,\ldots,\lambda_n$ (counted with respect to algebraic multiplicity), and all of these eigenvalues are real. Furthermore, there is an orthonormal basis $\{\bm{e}_1,\ldots,\bm{e}_n\}$ for $\mathbb{R}^n$, such that $\bm{e}_i$ is an eigenvector corresponding to $\lambda_i$.
\end{Theorem}

The proof of this is given in any good elementary linear algebra text. The result will be a corollary to the theorems we prove below about self-adjoint compact operators.

One of our first tasks is to generalize the concept of eigenvalues and eigenvectors to accommodate the operators considered in this section (which may be defined in infinite-dimensional spaces and may be unbounded).

\begin{Definition}
\label{def:shifted-operator-and-resolvent}
Let $X$ be a complex Banach space. Let $(\mathcal{D}(A),A)$ be an operator from $X$ to $X$. For any $\lambda\in\mathbb{C}$ we define the operator $(\mathcal{D}(A),A_\lambda)$ by
\begin{equation}\label{eq:shifted-operator-definition}
A_\lambda:=A-\lambda I,
\end{equation}
where $I$ is the identity operator on $X$. If $A_\lambda$ has an inverse (i.e., if it is one-to-one), we denote the inverse by $R_\lambda(A)$, and call it the \textbf{resolvent} of $A$.
\end{Definition}

\begin{Definition}
\label{def:resolvent-spectrum-and-spectral-parts}
Let $X\ne\{0\}$ be a complex Banach space and let $(\mathcal{D}(A),A)$ be a linear operator from $X$ to $X$. Consider the following three conditions:\begin{enumerate}\item $R_\lambda(A)$ exists,\item $R_\lambda(A)$ is bounded,\item the domain of $R_\lambda(A)$ is dense in $X$.\end{enumerate}We decompose the complex plane $\mathbb{C}$ into the following two sets.\begin{itemize}\item The \textbf{resolvent set} of the operator $A$ is the set
\begin{equation}\label{eq:resolvent-set-definition}
\rho(A):=\{\lambda\in\mathbb{C}\mid\text{(1), (2), and (3) hold}\}.
\end{equation}
Elements $\lambda\in\rho(A)$ in the resolvent set are called \textbf{regular values} of the operator $A$.\item The \textbf{spectrum} of the operator $A$ is the complement of the resolvent set
\begin{equation}\label{eq:operator-spectrum-definition}
\sigma(A):=\mathbb{C}\backslash\rho(A).
\end{equation}
The spectrum can be further decomposed into three disjoint sets.\begin{itemize}\item The \textbf{point spectrum} or \textbf{discrete spectrum} is the set
\begin{equation}\label{eq:operator-point-spectrum-definition}
\sigma_p(A):=\{\lambda\in\sigma(A)\mid\text{(1) does not hold}\}.
\end{equation}
That is, the point spectrum is the set of $\lambda\in\mathbb{C}$ for which $\mathcal{N}(A_\lambda)$ is nontrivial. Elements of the point spectrum are called \textbf{eigenvalues}. If $\lambda\in\sigma_p(A)$, elements $x\in\mathcal{N}(A_\lambda)$ are called \textbf{eigenvectors} or \textbf{eigenfunctions} of $A$. The dimension of $\mathcal{N}(A_\lambda)$ is called the (geometric) \textbf{multiplicity} of $\lambda$.\item The \textbf{continuous spectrum} is the set
\begin{equation}\label{eq:operator-continuous-spectrum-definition}
\sigma_c(A):=\{\lambda\in\sigma(A)\mid\text{(1) and (3) hold but (2) does not}\}.
\end{equation}
\item The \textbf{residual spectrum} or \textbf{compression spectrum} is the set
\begin{equation}\label{eq:operator-residual-spectrum-definition}
\sigma_r(A):=\{\lambda\in\sigma(A)\mid\text{(1) holds but (3) does not}\}.
\end{equation}
Since $\overline{\mathcal{R}(A_\lambda)}\ne X$ we say that the range has been compressed.\end{itemize}\end{itemize}
\end{Definition}

\begin{Definition}
\label{def:spectral-deficiency}
If $X$ is a Hilbert space, we refer to the dimension of $\mathcal{R}(A_\lambda)^\perp$ as the \textbf{deficiency} of $\lambda\in\mathbb{C}$.
\end{Definition}

Note that by our definition, $\lambda\in\sigma(A)$ can have nonzero deficiency and not be in the compression spectrum. Some authors define the compression spectrum to be all $\lambda\in\mathbb{C}$ such that the deficiency is nonzero, but in this case the point spectrum and compression spectrum are not necessarily disjoint.

\begin{Example}
\label{ex:finite-dimensional-spectrum-regular-eigenvalues}
One of the fundamental results of linear algebra is that for a linear operator $A$ on a finite-dimensional space the continuous spectrum and the compression spectrum of the operator are empty; i.e., the complex plane can be decomposed into regular values and eigenvalues of the operator.
\end{Example}

\begin{Example}
\label{ex:right-shift-compression-spectrum}
For a simple example of an operator with a spectral value that is not an eigenvalue, consider the right-shift operator $S_r:\ell^2\to\ell^2$. The complex number $\lambda=0$ is an element of the spectrum. To see this we recall that the resolvent operator $R_0(S_r)$ is simply the left-shift operator $S_l$ operating on the domain $\{1,0,0,\ldots\}^\perp$, and while this operator is bounded, its domain is not dense in $\ell^2$. Thus, $\lambda=0$ is in the compression spectrum of $S_r$ and has deficiency $1$.
\end{Example}

Spectral theory is a very broad and well studied subject. Our treatment of it here is of necessity very cursory; our aim is primarily to develop the tool of eigenfunction expansions. Thus, we begin with a basic theorem about eigenvectors.

\begin{Theorem}
\label{thm:distinct-eigenvalue-eigenvectors-independence}
If $\lambda_i$, $i=1,\ldots,n$, are distinct eigenvalues of the operator $(\mathcal{D}(A),A)$ and $x_i\in\mathcal{N}(A_{\lambda_i})$ are corresponding eigenvectors, then the set\[\{x_1,x_2,\ldots,x_n\}\]is linearly independent.
\end{Theorem}

\begingroup
\postdisplaypenalty=10000
\renewcommand{\endProof}{\par\nobreak\hfill$\blacksquare$\par\vspace{0.1cm}}
\begin{Proof}
Suppose not. Then there is an integer $k\in[2,n]$ such that the set $\{x_1,\ldots,x_{k-1}\}$ is linearly independent, whereas $x_k$ can be expanded in this set; i.e.,
\begin{equation}\label{eq:eigenvector-first-dependence-expansion}
x_k=\alpha_1x_1+\alpha_2x_2+\cdots+\alpha_{k-1}x_{k-1},
\end{equation}
where the coefficients $\alpha_i$ are not all zero. We now apply $(A-\lambda_kI)$ to both sides of the equation to get\[\begin{aligned}
0&=(A-\lambda_kI)x_k\\
&=(A-\lambda_kI)[\alpha_1x_1+\alpha_2x_2+\cdots+\alpha_{k-1}x_{k-1}]\\
&=\alpha_1(\lambda_1-\lambda_k)x_1+\alpha_2(\lambda_2-\lambda_k)x_2\\
&\quad+\cdots+\alpha_{k-1}(\lambda_{k-1}-\lambda_k)x_{k-1}.
\end{aligned}\]Since $\{x_1,\ldots,x_{k-1}\}$ is linearly independent we have
\begin{equation}\label{eq:eigenvector-independence-coefficient-identity}
(\lambda_i-\lambda_k)\alpha_i=0,\quad i=1,\ldots,k-1.
\end{equation}
However, since $\lambda_i\ne\lambda_m$ this implies $\alpha_i=0$, $i=1,\ldots,k-1$. This is a contradiction and completes our proof.
\end{Proof}
\endgroup
